
\section{Conclusion}\label{sec:conclusion}

This paper explained the design and implementation of the distributed data parallel module in PyTorch v1.5, and conducted performance evaluations on NCCL and Gloo backend using ResNet50 and BERT models. \code{DDP} accelerates training by aggregating gradients into buckets for communication, overlapping communication with computation, and skipping synchronizations. We also highlighted real-world caveats in gradient synchronization which are important for broad adoption.  Results showed that \code{DDP} with NCCL backend can achieve near-linear scalability on 256 GPUs when configured properly. The measurements also revealed that the backward pass in \code{DDP} is the most expensive step in training and requires efforts from both framework developers to enable optimization algorithms and application developers to empirically configure the knobs. Based on our observations, we shared lessons learned from serving a variety of application, discussed potential future improvements for distributed data parallel training, and enthusiastically encourage open source community to experiment with more novel ideas.

\newpage